\subsection{AMALGAMATED SUM OF MONOIDS}

Let \((M_{i})_{i\in I}\) denote a family of monoids and \(e_i\) the identity element of \(\mathbf{M}_i\) . We are given a monoid \(\mathbf{A}\) and a family of homomorphisms \(h_i\colon \mathbf{A}\to \mathbf{M}_i\) (for \(i\in I\) ).

The set S the sum of the family \((\mathbf{M}_{i})_{i\in\mathbb{I}}\) has elements the ordered pairs \((i,x)\) with \(i\in I\) and \(x\in M_{i}\) . For every triple \(\alpha=(i,x,x^{\prime})\) with \(i\in I\) , \(x,x^{\prime}\) in \(M_{i}\) , write \(u_{\alpha}=(i,xx^{\prime})\) and \(v_{\alpha}=(i,x).(i,x^{\prime})\) ; for every triple \(\lambda=(i,j,a)\) in \(I\times I\times A\) , write \(p_{\lambda}=(i,h_{i}(a))\) and \(q_{\lambda}=(j,h_{j}(a))\) ; for all \(i\in I\) , write \(\varepsilon_{i}=(i,e_{i})\) . The monoid M defined by S and the relators \((u_{\alpha},v_{\alpha})\) , \((p_{\lambda},q_{\lambda})\) and \((\varepsilon_{i},e)\) is called the sum of the family \((\mathbf{M}_{i})_{i\in\mathbb{I}}\) amalgamated by A. Let \(\phi\) denote the canonical homomorphism of \(\mathrm{Mo}(\mathrm{S})\) onto M and write \(\phi_{i}(x)=\phi(i,x)\) for \((i,x)\in\mathrm{S}\) . It is said that \(\phi_{i}\) is the canonical mapping of \(M_{i}\) into M. For all \(a\in A\) , the element \(\phi(i,h_{i}(a))\) is independent of i and denoted by \(h(a).†\)

The universal property of monoids defined by generators and relators (no. 2) implies the following result:

PROPOSITION 4. (a) For all \(i \in \mathbf{I}\) , the mapping \(\phi_i\) is a homomorphism of \(\mathbf{M}_i\) into \(\mathbf{M}\) and \(\phi_i \circ h_i = h\) for all \(i \in \mathbf{I}\) . Further, \(\mathbf{M}\) is generated by \(\bigcup_{i \in \mathbf{I}} \phi_i(\mathbf{M}_i)\) .

\begin{enumerate}
\def\labelenumi{(\alph{enumi})}
\setcounter{enumi}{1}
\tightlist
\item
  Let \(\mathbf{M}'\) be a monoid and \(f': \mathbf{M}_i \to \mathbf{M}'\) (for \(i \in \mathbf{I}\) ) homomorphisms such that \(f_i \circ h_i\) is independent of \(i \in \mathbf{I}\) . There exists one and only one homomorphism \(f: \mathbf{M} \to \mathbf{M}'\) such that \(f_i = f \circ \phi_i\) for all \(i \in \mathbf{I}\) .
\end{enumerate}

In what follows we shall make the following hypothesis:

\begin{enumerate}
\def\labelenumi{(\Alph{enumi})}
\tightlist
\item
  For all \(i \in \mathbf{I}\) , there exists a subset \(\mathbf{P}_i\) of \(\mathbf{M}_i\) containing \(e_i\) such that the mapping \((a, p) \mapsto h_i(a).p\) of \(\mathbf{A} \times \mathbf{P}_i\) into \(\mathbf{M}_i\) is bijective.
\end{enumerate}

It implies that the homomorphisms \(h_i\) are injective. Let \(x \in \mathbf{M}\) ; every finite sequence \(\sigma = (a: i_1, \ldots, i_n; p_1, \ldots, p_n)\) with \(a \in \mathbf{A}\) , \(i_\alpha \in \mathbf{I}\) and \(p_\alpha \in \mathbf{P}_{i_\alpha}\) for \(1 \leqslant \alpha \leqslant n\) , satisfying

\[
x = h (a) \cdot \prod_ {\alpha = 1} ^ {n} \phi_ {i _ {\alpha}} (p _ {\alpha})\tag{3}
\]

is called a decomposition of \(x\) . The integer \(n \geqslant 0\) is called the length of the decomposition \(\sigma\) and is denoted by \(l(\sigma)\) ; the sequence \((e)\) is a decomposition of length 0 of the identity element of \(M\) . The decomposition \(\sigma\) is called reduced if \(i_{\alpha} \neq i_{\alpha + 1}\) for \(1 \leqslant \alpha < n\) and \(p_{\alpha} \neq e_{i_{\alpha}}\) for \(1 \leqslant \alpha \leqslant n\) .

PROPOSITION 5. Under hypothesis (A) every element \(x\) of \(\mathbf{M}\) admits a unique reduced decomposition \(\sigma\) . Every decomposition \(\sigma' \neq \sigma\) of \(x\) satisfies \(l(\sigma') > l(\sigma)\) .

\BourbakiRunInHeading{(A) Uniqueness of a reduced decomposition:}

Let \(\Sigma\) denote the set of sequences \(\sigma = (a; i_1, \ldots, i_n; p_1, \ldots, p_n)\) with \(n \geqslant 0\) , \(a \in \mathbf{A}\) , \(i_\alpha \in \mathbf{I}\) and \(p_\alpha \in \mathrm{P}_{i_\alpha} - \{e_{i_\alpha}\}\) for \(1 \leqslant \alpha \leqslant n\) , such that \(i_\alpha \neq i_{\alpha+1}\) for \(1 \leqslant \alpha < n\) . Let \(\Phi\) denote the mapping of \(\Sigma\) into \(\mathbf{M}\) defined by

\[
\Phi (a; i _ {1}, \dots , i _ {n}; p _ {1}, \dots , p _ {n}) = h (a). \prod_ {\alpha = 1} ^ {n} \phi_ {i _ {\alpha}} (p _ {\alpha}).\tag{4}
\]

A reduced decomposition of \(x \in \mathbf{M}\) is an element \(\sigma\) of \(\Sigma\) such that \(\Phi(\sigma) = x\) . For all \(i \in \mathbf{I}\) , let \(\Sigma_i\) be the subset of \(\Sigma\) consisting of the sequences \((e; i_1, \ldots, i_n; p_1, \ldots, p_n)\) with \(i \neq i_1\) when \(n > 0\) . Let

\[
\sigma = (e; i _ {1}, \dots , i _ {n}; p _ {1}, \dots , p _ {n})
\]

be in \(\Sigma_{i}\) and \(\xi\) in \(\mathbf{M}_i\) ; let \(\xi = h_i(a).p\) with \(a\in \mathbf{A}\) and \(p\in \mathbf{P}_i\) , and

\[
\Psi_ {i} (\xi , \sigma) = \left\{ \begin{array}{l l} (a; i _ {1}, \dots , i _ {n}; p _ {1}, \dots , p _ {n}) & \text { if } p = e _ {i} \\ (a; i, i _ {1}, \dots , i _ {n}; p, p _ {1}, \dots , p _ {n}) & \text { if } p \neq e _ {i}. \end{array} \right.\tag{5}
\]

It is immediate that \(\Psi_{i}\) is a bijection of \(\mathbf{M}_i\times \Sigma_i\) onto \(\Sigma\) .

Let \(i \in I\) and \(x \in M_i\) ; as \(\Psi_i\) is bijective, a mapping \(f_{i,x}\) of \(\Sigma\) into itself is defined by

\[
f _ {i, x} (\Psi_ {i} (\xi , \sigma)) = \Psi_ {i} (x \xi , \sigma) \quad (\xi \in \mathbf {M} _ {i}, \sigma \in \Sigma_ {i}).\tag{6}
\]

Further, for \(a \in \mathbf{A}, f_a\) denotes the mapping of \(\Sigma\) into itself defined by

\[
f _ {a} (a ^ {\prime}; i _ {1}, \dots , i _ {n}; p _ {1}, \dots , p _ {n}) = (a a ^ {\prime}; i _ {1}, \dots , i _ {n}; p _ {1}, \dots , p _ {n}).\tag{7}
\]

Clearly \(f_{i, e_i}\) is the identity mapping of \(\Sigma\) and \(f_{i, xx'} = f_{i, x} \circ f_{i, x'}\) for \(x, x'\) in \(M_i\) and \(f_{i, h_i(a)} = f_a\) for \(a \in A\) and \(i \in I\) .

Then Proposition 4 may be applied to the case where \(\mathbf{M}'\) is the monoid of mappings of \(\Sigma\) into itself with law of composition \((f, f') \mapsto f \circ f'\) and where \(f_{i}\) is the homomorphism \(x \mapsto f_{i,x}\) of \(\mathbf{M}_{i}\) into \(\mathbf{M}'\) ; then there exists a homomorphism \(f\) of \(\mathbf{M}\) into \(\mathbf{M}'\) such that \(f_{i,x} = f(\phi_i(x))\) for \(i \in \mathbf{I}\) and \(x \in \mathbf{M}_i\) . Let

\[
\sigma = (a; i _ {1}, \dots , i _ {n}; p _ {1}, \dots , p _ {n})
\]

be in \(\Sigma\) . Formulae (5) to (7) imply by induction on \(n\) the relation

\[
\begin{array}{l} \sigma = (f _ {a} \circ f _ {i _ {1}, p _ {1}} \circ \dots \circ f _ {i _ {n}, p _ {n}}) (e) \\ = f (h (a) \phi_ {i _ {1}} (p _ {1}) \dots \phi_ {i _ {n}} (p _ {n})) (e), \end{array}
\]

that is \(\sigma = f(\Phi(\sigma))(e)\) . This proves that \(\Phi\) is injective.

\BourbakiRunInHeading{(B) Existence of a decomposition:}

Let D be the set of elements of M admitting a decomposition. Then \(e \in D\) and M is generated by \(\bigcup_{i \in I} \phi_i(M_i)\) and hence by \(h(A) \cup \bigcup_{i \in I} \phi_i(P_i)\) . Then \(D \cdot \phi_i(P_i) \subset D\) for all \(i \in I\) ; to prove that D = M, it thus suffices to prove the relation \(D \cdot h(A) \subset D\) . This follows from the following more precise lemma:

Lemma 1. Let \(i_1, \ldots, i_n\) be in I and \(p_\alpha\) in \(\mathbf{P}_{i_\alpha}\) for \(1 \leqslant \alpha \leqslant n\) . For all \(a \in A\) there exists \(a' \in A\) and a sequence \((p'_\alpha)_{1 \leqslant \alpha \leqslant n}\) with \(p'_\alpha \in \mathbf{P}_{i_\alpha}\) such that

\[
\phi_ {i _ {1}} (p _ {1}) \dots \phi_ {i _ {n}} (p _ {n}) h (a) = h (a ^ {\prime}) \phi_ {i _ {1}} (p _ {1} ^ {\prime}) \dots \phi_ {i _ {n}} (p _ {n} ^ {\prime}).
\]

\(h(a) = \phi_{i_n}(h_{i_n}(a))\) and there exists \(a_{n} \in \mathbf{A}\) and \(p_n' \in \mathbf{P}_{i_n}\) with

\[
p _ {n} \cdot h _ {i _ {n}} (a) = h _ {i _ {n}} (a _ {n}) \cdot p _ {n} ^ {\prime}.
\]

It follows that \(\phi_{i_n}(p_n)h(a) = h(a_n)\phi_{i_n}(p_n')\) , whence

\[
\phi_ {i _ {1}} (p _ {1}) \dots \phi_ {i _ {n - 1}} (p _ {n - 1}) \phi_ {i _ {n}} (p _ {n}) h (a) = \phi_ {i _ {1}} (p _ {1}) \dots \phi_ {i _ {n - 1}} (p _ {n - 1}) h (a _ {n}) \phi_ {i _ {n}} (p _ {n} ^ {\prime});
\]

the lemma follows from this by induction on n.

\BourbakiRunInHeading{(C) End of the proof:}

Let \(x \in M\) and let n be the minimum of the lengths of decompositions of x. We shall prove that every decomposition \(\sigma\) of x of length n is reduced. This will establish the existence of a reduced decomposition of x; the uniqueness of the reduced decomposition then implies \(l(\sigma') > l(\sigma)\) for every decomposition \(\sigma' \neq \sigma\) of x.

The case \(n = 0\) being trivial, suppose \(n > 0\) . Let

\[
\sigma = (a; i _ {1}, \dots , i _ {n}; p _ {1}, \dots , p _ {n})
\]

be a decomposition of \(x\) of length \(n\) . If there existed an integer \(\alpha\) with \(1 \leqslant \alpha \leqslant n\) and \(p_{\alpha} = e_{i_{\alpha}}\) , the sequence

\[
(a; i _ {1}, \dots , i _ {\alpha - 1}, i _ {\alpha + 1}, \dots , i _ {n}; p _ {1}, \dots , p _ {\alpha - 1}, p _ {\alpha + 1}, \dots , p _ {n})
\]

would be a decomposition of \(x\) of length \(n - 1\) , which is excluded. Suppose that there exists an integer \(\alpha\) with \(1 \leqslant \alpha < n\) and \(i_{\alpha} = i_{\alpha + 1}\) and let

\[
p _ {\alpha} p _ {\alpha + 1} = h _ {i _ {\alpha}} \left(a ^ {\prime}\right) \cdot p _ {\alpha} ^ {\prime}
\]

with \(a' \in \mathbf{A}\) and \(p_{\alpha}' \in \mathrm{P}_{t_{\alpha}}\) ; by Lemma 1 there exists elements \(a'' \in \mathbf{A}\) ,

\[
p _ {1} ^ {\prime} \in \mathrm{P} _ {i _ {1}}, \dots , p _ {\alpha - 1} ^ {\prime} \in \mathrm{P} _ {i _ {\alpha - 1}}
\]

such that

\[
\phi_ {i _ {1}} (p _ {1}) \dots \phi_ {i _ {\alpha - 1}} (p _ {\alpha - 1}) h (a ^ {\prime}) = h (a ^ {\prime \prime}) \phi_ {i _ {1}} (p _ {1} ^ {\prime}) \dots \phi_ {i _ {\alpha - 1}} (p _ {\alpha - 1} ^ {\prime})
\]

and the sequence

\[
\left(a a ^ {\prime \prime}; i _ {1}, \dots , i _ {\alpha - 1}, i _ {\alpha}, i _ {\alpha + 2}, \dots , i _ {n}; p _ {1} ^ {\prime}, \dots , p _ {\alpha - 1} ^ {\prime}, p _ {\alpha} ^ {\prime}, p _ {\alpha + 2}, \dots , p _ {n}\right)
\]

is a decomposition of \(x\) of length \(n - 1\) , which is a contradiction.

We have thus proved that \(\sigma\) is reduced.

COROLLARY. Under hypothesis (A) the homomorphisms \(\phi_i\) and \(h\) are injective. For \(i \neq j\) in I, \(\phi_i(M_i) \cap \phi_j(M_j) = h(A)\) .

First \(h\) is injective: if \(h(a) = h(a')\) , then (a) and \((a')\) are two reduced decompositions of the same element of \(\mathbf{M}\) , whence \(a = a'\) . Let \(i \in \mathbf{I}\) ; then \(h(\mathbf{A}) = \phi_i(h_i(\mathbf{A})) \subset \phi_i(\mathbf{M}_i)\) ; the uniqueness of reduced decompositions implies

\[
h (\mathrm{A}) \cap \phi_ {i} (\mathrm{M} _ {i} - h _ {i} (\mathrm{A})) = \varnothing ,
\]

whence \(\phi_{i}(\mathbf{M}_{i} - h_{i}(\mathbf{A})) = \phi_{i}(\mathbf{M}_{i}) - h(\mathbf{A}).\)

The injectivity of the homomorphisms \(\phi_{i}\) and the relation

\[
\phi_ {i} (\mathbf {M} _ {i}) \cap \phi_ {j} (\mathbf {M} _ {j}) \subset h (\mathrm{A})
\]

for \(i \neq j\) are then consequences of the following fact: for \(i, j\) in \(I, x\) in \(M_i = h_i(A)\) and \(y\) in \(M_j = h_j(A)\) , the relation \(\phi_i(x) = \phi_j(y)\) implies \(i = j\) and \(x = y\) . Let \(x = h_i(a).p\) and \(y = h_j(b).q\) with \(a, b\) in \(A, p\) in \(P_i = \{e_i\}\) and \(y\) in \(P_j = \{e_j\}\) . Then \(\phi_i(x) = h(a)\phi_i(p)\) and \(\phi_j(y) = h(b)\phi_j(q)\) and hence \((a; i; p)\) and \((b; j; q)\) are two reduced decompositions of the same element of \(M\) . It follows that \(i = j, a = b\) and \(p = q\) , whence \(x = h_i(a)p = h_j(b)q = y\) .

When hypothesis (A) is fulfilled, we shall identify each monoid \(M_{i}\) with a submonoid of M by means of \(\phi_{i}\) ; similarly, we shall identify A with a submonoid of M by h. Then M is generated by \(\bigcup_{i\in I}M_{i}\) and \(M_{i}\cap M_{j}=A\) for \(i\neq j\) .

Every element of M can be written uniquely in the form \(a.\prod_{\alpha=1}p_{\alpha}\) with \(a\in A, p_{1}\in P_{i1}=\{e\},\ldots,p_{n}\in P_{in}=\{e\}\) and \(i_{\alpha}\neq i_{\alpha+1}\) for \(1\leqslant\alpha<n\) . Finally, if \(M'\) is a monoid and \((f_{i}:M_{i}\to M')\) (for \(i\in I\) ) a family of homomorphisms whose restrictions to A are the same homomorphism of A into \(M'\) , there exists one and only one homomorphism \(f:M\to M'\) inducing \(f_{i}\) on \(M_{i}\) for all \(i\in I\) .

Hypothesis (A) is satisfied in two important cases:

\begin{enumerate}
\def\labelenumi{(\alph{enumi})}
\tightlist
\item
  \(\mathbf{A} = \{e\}\) . In this case, there is a family \((\mathbf{M}_i)_{i\in I}\) of monoids and \(\mathbf{M}\) is called the monoidal sum of this family. Each \(\mathbf{M}_i\) is identified with a submonoid of \(\mathbf{M}\) and \(\mathbf{M}\) is generated by \(\bigcup_{i\in I}\mathbf{M}_i\) ; further, \(\mathbf{M}_i\cap \mathbf{M}_j = \{e\}\) for \(i\neq j\) . Every element of \(\mathbf{M}\) may be written uniquely in the form \(x_{1}\ldots x_{n}\) with
\end{enumerate}

\[
x _ {1} \in \mathrm{M} _ {i _ {1}} - \{e \}, \dots , x _ {n} \in \mathrm{M} _ {i _ {n}} - \{e \}
\]

and \(i_{\alpha} \neq i_{\alpha + 1}\) for \(1 \leqslant \alpha \leqslant n\) . Finally, for every family of homomorphisms \((f_i: \mathbf{M}_i \to \mathbf{M}')\) , there exists a unique homomorphism \(f: \mathbf{M} \to \mathbf{M}'\) whose restriction to \(\mathbf{M}_i\) is \(f_i\) for all \(i \in I\) .

\begin{enumerate}
\def\labelenumi{(\alph{enumi})}
\setcounter{enumi}{1}
\tightlist
\item
  There is a family of groups \((\mathbf{G}_i)_{i\in I}\) containing as subgroup the same group A and \(h_i\) is the injection of A into \(\mathbf{G}_i\) . The sum of the family \((\mathbf{G}_i)_{i\in I}\) amalgamated by A is then a group G: the monoid G is generated by \(\bigcup_{i\in I}\phi_i(\mathbf{G}_i)\) and every element of \(\bigcup_{i\in I}\phi_i(\mathbf{G}_i)\) admits an inverse in G (cf.~§ 2, no. 3, Corollary 1 to Proposition 4); it is denoted by \(\star_{\mathrm{A}}\mathrm{G}_{i}\) or \(\mathrm{G}_{1}*\mathrm{A}\mathrm{G}_{2}\) when \(\mathbf{I} = \{1,2\}\) . When A consists of the identity element, it is also said that G is the free product of the family \((\mathbf{G}_i)_{i\in I}\) of groups and it is denoted by \(\star \mathbf{G}_i\) (or \(\mathbf{G}_1*\mathbf{G}_2\) if \(\mathbf{I} = \{1,2\}\) ).†
\end{enumerate}

